\subsection{Fibrations with positive-genus ruling base}
\label{subsec:positive-genus-picard-fibration}

When the ruling curve of the distinguished normalized cover has positive
genus, its Frobenius twist is a base of a fibration on the smooth quotient.
The resulting morphism on the normalization controls its singularities
and the arithmetic genus of the quotient fibres. We first record a
consequence of the divided differential that applies to every conductor
equality.

\begin{proposition}\label{prop:equality-rational-curves}
Let $X$ be a smooth projective integral surface over an algebraically
closed field $k$ of characteristic $p>0$. Let $Y\subset X$ be a smooth
integral ample divisor, and let $M\ne\OO_X$ have order $p$ with
$M|_Y\simeq\OO_Y$. Suppose that
\[
 2g(Y)-2=apY^2
\]
for the conductor integer $a\geq1$ of
Proposition~\ref{prop:normalized-mup-conductor}.
Every smooth rational curve $E\subset X$ satisfies
\[
 E^2\leq-pa(Y\cdot E).
\]
In particular $X$ has no $(-1)$-curve.
\end{proposition}

\begin{proof}
Put $m=pa$. Proposition~\ref{prop:normalized-mup-singularity-defect}
gives a nonzero section
\[
 \theta\in H^0\bigl(X,\Omega^1_{X/k}(-mY)\bigr)
\]
whose zero scheme is finite. Let $e=Y\cdot E$, which is a positive
integer by ampleness. The twisted conormal sequence of $E$ is
\[
 0\longrightarrow\OO_E(-E-mY)\longrightarrow
 \Omega^1_{X/k}(-mY)|_E\longrightarrow
 \Omega^1_{E/k}\otimes\OO_E(-mY)\longrightarrow0.
\]
Its last term is $\OO_{\PP^1}(-2-me)$ and has no nonzero section.
Consequently $\theta|_E$ comes from a section of
$\OO_E(-E-mY)$. That section is nonzero, since otherwise the zero
scheme of $\theta$ would contain $E$. Its degree is therefore
nonnegative, giving $-E^2-me\geq0$. If $E^2=-1$, this would imply
$-1\leq-pae\leq-2$, a contradiction.
\end{proof}

We will use a genus bound that holds for singular integral curves as
well as smooth ones.

\begin{lemma}\label{lem:very-ample-curve-genus-bound}
Let $F$ be an integral projective curve over an algebraically closed
field, and let $L$ be a very ample invertible sheaf of degree $b$ on $F$.
Then
\[
 p_a(F)\leq\frac{(b-1)(b-2)}2.
\]
\end{lemma}

\begin{proof}
Embed $F$ by $L$, and let $R$ be its standard graded homogeneous
coordinate ring. Choose a hyperplane equation $\ell$ that does not
vanish identically on $F$. Multiplication by $\ell$ is injective because
$R$ is a domain. Put $Q=R/(\ell)$, and let $T$ be the saturated
homogeneous coordinate ring of the finite hyperplane section
$\Gamma=F\cap V(\ell)$, whose length is $b$. There is a graded
surjection $Q\to T$. Writing $H$ for dimensions of graded pieces, we get
\[
 H_R(n)-H_R(n-1)=H_Q(n)\geq H_T(n)\qquad(n\geq1).
\]

Choose a linear form $z$ not vanishing at any point of $\Gamma$.
It is a nonzerodivisor on $T$: its kernel is zero after sheafification
on $\Gamma$, and saturation removes the irrelevant-ideal torsion.
Hence
\[
 H_T(n)-H_T(n-1)=\dim_k(T/zT)_n\geq0.
\]
The algebra $T/zT$ is standard graded. If one of its graded pieces is
zero, all later pieces are zero, since the algebra is generated in
degree one. Thus $H_T$ increases strictly until it reaches its eventual
value $b$. Since $H_T(0)=1$, this proves
\[
 H_T(n)\geq\min(n+1,b)\qquad(n\geq0).
\]
For $n$ sufficiently large the Hilbert polynomial of $F$ gives
$H_R(n)=bn+1-p_a(F)$, whereas $H_R(0)=1$. Summing the first inequality
yields
\[
 \begin{split}
 bn+1-p_a(F)
 &\geq1+\sum_{j=1}^{n}\min(j+1,b)\\
 &=bn+1-\frac{(b-1)(b-2)}2,
 \end{split}
\]
which is the required bound.
\end{proof}

\begin{proposition}\label{prop:positive-genus-picard-fibration}
Let $X$ be a smooth projective integral surface over an algebraically
closed field $k$ of characteristic $p>0$, let $Y\subset X$ be a smooth
integral ample divisor, and let $M\ne\OO_X$ have order $p$ with
$M|_Y\simeq\OO_Y$. Suppose that
\[
 2g(Y)-2=pY^2.
\]
Let $\pi:N\to X$ be the normalization of the $\mu_p$-torsor defined
by a trivialization of $M^p$, and put $D=\pi^*Y$.
Assume that $N$ is birationally ruled, and fix a presentation
\[
 k(N)=k(C)(t)
\]
with $C$ a smooth projective integral curve of genus at least one.
Put $B=C^{(p)}$, and denote relative Frobenius by
$F_{C/k}:C\to B$.

There are projective flat morphisms $f:X\to B$ and $h:N\to C$ such that
\[
 f\pi=F_{C/k}h,
\]
and $N$ is the normalization of the integral surface $X\times_B C$.
The singularities of $N$ are rational double points. Moreover,
\[
 h_*\OO_N=\OO_C,\qquad R^ih_*\OO_N=0\quad(i>0),
 \qquad \chi(\OO_N)=1-g(C).
\]

There is a nonempty open subset $C_0\subset C$ over which
$N_{C_0}\simeq C_0\times\PP^1$ as a $C_0$-scheme.
For corresponding general closed points
$c\in C$ and $b_0=F_{C/k}(c)\in B$, let
$F_N=h^{-1}(c)$ and $F_X=f^{-1}(b_0)$. Then $F_X$ is integral and
geometrically unibranch, and
$\pi|_{F_N}:F_N\simeq\PP^1\to F_X$ is its finite radicial
normalization. One has
\[
 \pi^*F_X=pF_N,\qquad \pi_*F_N=F_X.
\]
Writing $b=Y\cdot F_X=D\cdot F_N$, one obtains
\[
 2p_a(F_X)=(p-1)b,\qquad b>0.
\]
In characteristic two, $b$ is even. In particular the general fibre
$F_X$ is singular. If $Y$ is very ample, then $b\geq p+2$.
\end{proposition}

\begin{proof}
The equality conclusions of
Propositions~\ref{prop:normalized-mup-conductor} and
\ref{prop:normalized-mup-singularity-defect} give that $N$ is normal
and Gorenstein, $\pi$ is finite flat and radicial of degree $p$, and
\begin{equation}\label{eq:positive-genus-canonical}
 \omega_N\simeq\pi^*\omega_X\otimes\OO_N\bigl(-(p-1)D\bigr).
\end{equation}
Put $A=k(C)$, $K=k(X)$ and $L=k(N)=A(t)$.
A smooth ruled model of this assumed function field maps rationally
and purely inseparably to $X$ with degree $p$.
Thus Proposition~\ref{prop:distinguished-ruled-cover} applies and gives
\[
 A\cap K=A^p,\qquad AK=L,
\]
and a proper birational morphism $r:S\to N$, where $S$ is obtained
from $C\times\PP^1$ by finitely many point blowups. Write $q:S\to C$
for the resulting ruling morphism.

\emph{Extension of the map on the quotient.}
The inclusion $A^p\subset K$ defines a rational map $X\dashrightarrow B$.
Resolve it by point blowups $\rho:X'\to X$, with resulting morphism
$f':X'\to B$ \cite[Tag~0C5H]{Stacks}. These blowups remain regular,
and hence smooth over $k$ \cite[Tags~0AGP and~0B8X]{Stacks}.
Each nontrivial fibre of $\rho$ has connected reduced support whose
components are smooth rational curves. Indeed this holds for one point
blowup, and successive blowups insert projective lines meeting the
strict transforms of the components through the centre.

There is no nonconstant morphism $\PP^1\to B$. To check this also for
inseparable maps, suppose that such a map exists. If its map on
differentials is zero, its function-field image lies in $k(t^p)$;
replace $k(t)$ by that rational function field and repeat.
This process stops: a fixed nonconstant rational function has a zero
or pole of finite nonzero order and cannot belong to $k(t^{p^e})$
for every $e$. We obtain a map $\varphi:\PP^1\to B$ with nonzero
differential. The resulting nonzero map of line bundles
$\varphi^*\omega_B\to\omega_{\PP^1}$ would give an effective divisor
of degree
\[
 -2-\deg(\varphi)\bigl(2g(B)-2\bigr)<0,
\]
which is impossible because $g(B)=g(C)\geq1$.
Thus $f'$ is constant on each component of a $\rho$-fibre;
connectedness makes the values agree on the whole fibre.

Let $\Gamma$ be the reduced closed image of
$(\rho,f'):X'\to X\times B$. It is integral, and its projection to $X$
is proper and birational. Every fibre of this projection has one point
in its underlying set, because it is the image of a $\rho$-fibre.
The projection is therefore finite \cite[Tag~02LS]{Stacks}, and is an
isomorphism because $X$ is normal \cite[Tag~0AB1]{Stacks}.
Its second projection gives the required morphism $f:X\to B$.

\emph{The morphism on the normalization.}
The inclusion $A\subset L$ defines a rational lift to $C$.
Let $V=\Spec R\subset B$ be affine, and write
$F_{C/k}^{-1}(V)=\Spec R'$, with $R'$ finite over $R$.
Every element of $R'$ lies in $L$ and is integral over each local ring
of $N$ above $V$. It belongs to those local rings by normality, hence
is regular there. These maps glue to $h:N\to C$, and
$f\pi=F_{C/k}h$ because the equality holds at the generic point.
Likewise $hr=q$, since both morphisms to the separated curve $C$
agree generically.

Relative Frobenius on the smooth curve is finite flat of degree $p$:
its generic degree is $[A:A^p]=p$, and its finite local modules over
the DVRs of $B$ are torsion free. If $\alpha\in A\setminus A^p$, then
$A=A^p(\alpha)$. The polynomial $T^p-\alpha^p$ has no root in $K$,
since a root would be the unique element $\alpha$ and contradict
$A\cap K=A^p$. Thus
\[
 K\otimes_{A^p}A=K(\alpha)=L.
\]
The finite locally free algebra of $X\times_B C$ injects into this
generic field, so the fibre product is integral. Transitivity of
integrality shows that its normalization is the normalization of
$X$ in $L$, namely $N$.

Both $f$ and $h$ are projective, by the closed-graph construction from
the projective $k$-surfaces. They are flat: the local rings of their
integral dominant sources are torsion free over the local DVRs of
the base curves, and a torsion-free module over a DVR is flat.
Their closed fibres are Cartier divisors of pure dimension one.
Also
\begin{equation}\label{eq:positive-genus-functions}
 f_*\OO_X=\OO_B,\qquad h_*\OO_N=\OO_C.
\end{equation}
Indeed proper-cohomology finiteness makes these pushforwards finite
algebras on the base \cite[III, Theorem~3.2.1]{EGA}.
An element of $A(t)$ algebraic over $A$ belongs to $A$; consequently
the algebraic closure of $A^p$ in $K$ is $A\cap K=A^p$.
Thus the two finite algebras lie in their respective base function
fields and equal the base structure sheaves by normality.

\emph{Vanishing and rational singularities.}
For a point blowup $\beta$ of a smooth surface, one has
$\beta_*\OO=\OO$ and $R^i\beta_*\OO=0$ for $i>0$
\cite[Tag~0AGP, Lemma~54.3.4]{Stacks}.
The vanishing can be seen directly. For a regular local ring $R$
with parameters $x,y$, the two blowup charts have rings
$R[y/x]$ and $R[x/y]$, and their intersection has ring
$R[y/x,x/y]$. Every Laurent polynomial in $y/x$ is a sum of one
with nonnegative powers and one with negative powers. The degree-one
\v{C}ech differential is therefore surjective, and the cover has no
higher \v{C}ech term. The degree-zero pushforward follows from
normality and proper-cohomology finiteness.
The same two-chart calculation on $\PP^1_R$ gives $H^0=R$ and
vanishing higher cohomology. Successive applications of Leray yield
\[
 q_*\OO_S=\OO_C,\qquad R^iq_*\OO_S=0\quad(i>0).
\]

We also need $R^ih_*\OO_N=0$ for $i\geq2$, without any smoothness
assumption on the special fibres. Fix a closed point $c\in C$ and a
projective embedding $N\hookrightarrow\PP^m_C$.
Choose a hyperplane $H_1$ containing no component of $N_c$, and then
a hyperplane $H_2$ avoiding the finite scheme's support
$H_1\cap N_c$. These choices are possible over the infinite field $k$.
Lift their equations to a neighbourhood of $c$. The closed subset
$N\cap H_1\cap H_2$ has proper image in $C$ missing $c$.
Delete that image and shrink to an affine neighbourhood $V$.
Then $N_V\setminus H_1$ and $N_V\setminus H_2$ form an affine cover
of $N_V$, with affine intersection. The \v{C}ech complex gives
$H^i(N_V,\OO_N)=0$ for $i\geq2$. Such neighbourhoods prove the
claimed higher-direct-image vanishing.

Normality gives $r_*\OO_S=\OO_N$. The coherent sheaf
$\mathcal T=R^1r_*\OO_S$ is supported on finitely many points:
the proper birational morphism is an isomorphism over an open
containing every codimension-one point of $N$
\cite[Tag~0BFP]{Stacks}.
The low-degree Leray sequence for $q=hr$ is
\[
 0\longrightarrow R^1h_*\OO_N\longrightarrow R^1q_*\OO_S
 \longrightarrow h_*\mathcal T\longrightarrow R^2h_*\OO_N.
\]
The two established vanishings give $R^1h_*\OO_N=0$ and
$h_*\mathcal T=0$. The latter implies $\mathcal T=0$.
Indeed the closed subscheme cut out by
$\operatorname{Ann}(\mathcal T)$ is a zero-dimensional finite-type
$k$-scheme, hence is finite over $k$ and over $C$.
Pushforward along this finite morphism is faithful, being restriction
of scalars on affine opens.

Thus $N$ has rational singularities. To compare with the definition
requiring vanishing on every normal modification
\cite[Tag~0B4N]{Stacks}, localize at a singular closed point of $N$.
Regular points already have rational singularities.
The vanishing just proved gives $H^1(\OO)=0$ on the localized smooth
model. Any other normal modification is dominated by point blowups
of that model, by resolving the rational map to it with Tag~0C5H.
These blowups preserve the vanishing, and Leray injects the
$H^1(\OO)$ of the normal modification into that of the dominating
smooth model. It too is zero.

The local rings at singular points of $N$ are two-dimensional
normal Gorenstein algebraic local rings. They are Nagata and have
normal completions, by excellence
\cite[IV, 7.8.3(ii), (iii), (v), (vi)]{EGA}.
The classical rational-Gorenstein criterion therefore makes them
rational double points in every characteristic
\cite[Sections~54.9 and~54.12, Tags~0B4V and~0BGB]{Stacks}.
Finally, \eqref{eq:positive-genus-functions} and the vanishing of all
higher direct images of $h$ give
\[
 H^i(N,\OO_N)\simeq H^i(C,\OO_C)\quad(i\geq0),
\]
and hence $\chi(\OO_N)=1-g(C)$.

\emph{The general fibres and their multiplicities.}
Remove from $C$ the images of the finitely many points of $N$ over
which $r$ is not an isomorphism, and the images of all blowup centres
in $S\to C\times\PP^1$. Over the remaining nonempty open $C_0$,
both maps are isomorphisms, giving
$N_{C_0}\simeq C_0\times\PP^1$.

We check that the corresponding general fibre on $X$ is reduced.
Choose $v\in A^p\setminus(A^p)^p$. If $v$ were a $p$th power in $K$,
its unique root would belong to $A\cap K=A^p$, a contradiction.
Thus $dv\ne0$ in $\Omega^1_{K/k}$ \cite[Tag~031W]{Stacks}.
The map $f^*\Omega^1_{B/k}\to\Omega^1_{X/k}$ is consequently a
split injection on a dense open subset of $X$. The differential
criterion for morphisms between smooth $k$-schemes makes $f$ smooth
on this open \cite[IV, Theorem~17.11.1]{EGA}. Its image contains a
nonempty open of $B$.

Choose $b_0$ in that open and under $C_0$, and let $c$ be its unique
inverse image under Frobenius. Since $\pi$ is a universal
homeomorphism and $f\pi=F_{C/k}h$, the support of $F_X$ is identified
with the irreducible support of $F_N=\PP^1$.
Moreover $F_X$ meets the smooth locus of $f$, so it is generically
reduced. It is a Cartier divisor in a regular surface, hence is
Cohen--Macaulay and has no embedded associated points.
Generic reducedness then implies reducedness: a nonzero nilpotent
ideal would have an associated prime, necessarily a minimal prime,
contradicting its zero localization there. Thus $F_X$ is integral,
and it is geometrically integral over the algebraically closed field.

Pulling back the point $b_0$ in the commuting square gives the
Cartier-divisor identity
\[
 \pi^*F_X=h^*F_{C/k}^*[b_0]=pF_N.
\]
The residue degree along this divisor is one. Indeed, localizing at
the generic point of $F_X$ gives a finite free extension of DVRs of
rank $p$, with a unique maximal ideal upstairs and ramification
index $p$. Modulo the downstairs uniformizer, its dimension over
the downstairs residue field is both $p$ and $p$ times the residue
degree. Thus $\pi_*F_N=F_X$, and the finite morphism
$F_N\to F_X$ is birational. As $F_N$ is normal, it is the
normalization of $F_X$. The restriction remains radicial, proving
geometric unibranchness.

\emph{Arithmetic genus and the conditional degree bound.}
The projection formula gives
\[
 b=Y\cdot F_X=D\cdot F_N>0.
\]
This equality also follows directly on the singular curve $F_X$.
For its normalization $\nu$, the cokernel of
$\OO_{F_X}\to\nu_*\OO_{\PP^1}$ has finite length. Tensoring by an
invertible sheaf preserves that length; subtracting the two
Euler-characteristic equalities shows that normalization preserves
the degree of an invertible sheaf.

The curve $F_N$ lies in the regular locus of $N$, and $F_N^2=0$.
Adjunction therefore gives $K_N\cdot F_N=-2$.
Using \eqref{eq:positive-genus-canonical} and the residue degree one,
we obtain
\[
 -2=K_X\cdot F_X-(p-1)b.
\]
The integral Cartier fibre $F_X$ also has square zero. Adjunction
for this possibly singular curve gives
\[
 2p_a(F_X)-2=(K_X+F_X)\cdot F_X=K_X\cdot F_X;
\]
the Cartier adjunction formula applies without a smoothness
assumption on the curve \cite[Tags~0AA4 and~0A9X]{Stacks}.
Consequently $2p_a(F_X)=(p-1)b$. It is positive, whereas the
normalization has genus zero, so $F_X$ is singular.
When $p=2$, this genus formula also forces $b$ to be even.

Now suppose that $Y$ is very ample.
Its restriction $\OO_X(Y)|_{F_X}$ is very ample of degree $b$.
Lemma~\ref{lem:very-ample-curve-genus-bound} implies
\[
 (p-1)b\leq(b-1)(b-2),
 \qquad b^2-(p+2)b+2\geq0.
\]
The last quadratic is negative at both endpoints of the interval
$1\leq b\leq p+1$, and by convexity it is negative throughout that
interval. Hence $b\geq p+2$ under this additional hypothesis.
\end{proof}

\begin{remark}\label{rem:positive-genus-higher-conductor}
The fibration and rational-singularity arguments use the compatible
ruled model and the field equalities of
Proposition~\ref{prop:distinguished-ruled-cover}; they do not use
$a=1$ or very ampleness. Suppose more generally that $X,Y,M$ have
the same smoothness, projectivity and ampleness hypotheses, that
$a\geq1$ is the actual conductor integer of
Proposition~\ref{prop:normalized-mup-conductor}, and that
$2g(Y)-2=apY^2$. Assume that the distinguished normalization $N$
is birationally ruled with a positive-genus ruling base.
Then the same arguments apply, with
\[
 K_N=\pi^*K_X-a(p-1)D,
\]
and give
\[
 2p_a(F_X)=a(p-1)b.
\]
For $p=2$, the product $ab$ is even.
If $Y$ is in addition very ample, the embedded-curve genus bound gives
\[
 b\geq a(p-1)+3.
\]
The ruledness hypothesis is independent of the numerical equality.
It is also supplied if $X$ admits a degree-$p$ purely inseparable
ruled rational cover, by
Proposition~\ref{prop:distinguished-ruled-cover}.
\end{remark}

These conclusions apply to a merely ample divisor under the explicit
ruledness hypothesis. They do not identify $M$ with a pullback from
$B$, or make $h$ a projective-line bundle over all of $C$.
The rational-curve inequality in
Proposition~\ref{prop:equality-rational-curves} does not apply to the
singular general fibre $F_X$, even though its normalization is
$\PP^1$. The strict very ample theorem does not exclude these
merely ample configurations.
